\appendices
\section{Bộ Phân Tích CVE Dựa Trên Quy Tắc để Trích Xuất Syscall}
\label{appendix:cve_parser}

\subsection{Phạm Vi và Nguyên Tắc Thiết Kế}
\label{subsec:appendix_scope}

Chúng tôi thiết kế một bộ phân tích tất định, dựa trên quy tắc (thay vì NLP dựa trên ML) để trích xuất
các syscall có rủi ro cao từ mô tả CVE và khuyến cáo của nhà cung cấp, đảm bảo:

\begin{itemize}
  \item \textbf{Tính minh bạch:} Khớp mẫu hoàn toàn có thể kiểm tra và giải thích được.
  \item \textbf{Độ chính xác:} Sử dụng từ điển syscall đặc thù của kernel để tránh kết quả dương tính giả.
  \item \textbf{Nhận thức ngữ cảnh:} Lọc dựa trên điều kiện khai thác, phiên bản kernel,
        kiến trúc, các quyền hạn và ràng buộc thiết bị.
\end{itemize}

Phương pháp dựa trên quy tắc được ưu tiên hơn NLP dựa trên ML vì nó cung cấp tính tất định cao hơn,
không tạo ra kết quả dương tính giả từ suy luận mô hình, và cho phép kiểm tra đầy đủ từng
quyết định lọc. Điều này đặc biệt quan trọng trong các ngữ cảnh bảo mật, nơi khả năng kiểm tra
và giải thích là tối quan trọng.

\subsection{Xây Dựng Từ Điển Syscall}
\label{subsec:appendix_dictionary}

Từ điển syscall được xây dựng trực tiếp từ mã nguồn kernel:

\begin{lstlisting}[float=htbp, language=bash, caption={Trích xuất tên syscall từ mã nguồn kernel.}]
# Generated from kernel v5.15 headers
grep -r 'SYSCALL_DEFINE' arch/x86/kernel/ \
    | awk -F'(' '{print $2}' \
    | cut -d, -f1 \
    | sort -u
# --> [accept, access, acct, add_key, ...]
\end{lstlisting}

Nguồn tham khảo:
\begin{itemize}
  \item \texttt{/arch/x86/entry/syscalls/syscall\_64.tbl}
  \item \texttt{/include/linux/syscalls.h}
  \item Linux man-pages Phần~II
\end{itemize}

Điều này đảm bảo rằng chỉ các syscall hợp lệ trong kernel mới là ứng viên để trích xuất và phân tích. Từ
điển kết quả chứa tên chuẩn của 335~syscall được hỗ trợ bởi Linux kernel~6.x,
cùng với các bí danh thông dụng (ví dụ: \texttt{getpid} và \texttt{sys\_getpid}).

\subsection{Quy Tắc Trích Xuất và Lọc}
\label{subsec:appendix_extraction_rules}

Các mẫu regex được áp dụng cho mô tả CVE, khuyến cáo của nhà cung cấp và các
tiết lộ trên danh sách gửi thư bảo mật để xác định các đề cập syscall và lọc dựa trên ngữ cảnh.

\textbf{Các mẫu regex chính:}

\begin{lstlisting}[float=htbp, language=Python, caption={Các mẫu regex để xác định đề cập syscall trong mô tả CVE.}]
# Pattern 1: Direct syscall mention (e.g., "via the ptrace() syscall")
SYSCALL_DIRECT = re.compile(
    r'\b(?:the\s+)?'
    r'(?P<syscall>' + '|'.join(SYSCALL_DICT) + r')'
    r'(?:\(\d*\))?\s*(?:system\s+call|syscall)',
    re.IGNORECASE
)

# Pattern 2: Function name in exploit context
SYSCALL_EXPLOIT = re.compile(
    r'(?:exploit|abuse|invoke|call|use)\s+'
    r'(?:the\s+)?'
    r'(?P<syscall>' + '|'.join(SYSCALL_DICT) + r')'
    r'(?:\(\d*\))?',
    re.IGNORECASE
)

# Pattern 3: C annotation pattern (e.g., "setuid(0)" or "execve('/bin/sh', ...)")
SYSCALL_CODE = re.compile(
    r'\b(?P<syscall>' + '|'.join(SYSCALL_DICT) + r')'
    r'\s*\([^)]*\)',
    re.IGNORECASE
)
\end{lstlisting}

\textbf{Quy tắc lọc theo ngữ cảnh:}

Sau khi trích xuất các syscall ứng viên, bộ phân tích lọc chúng dựa trên:
\begin{enumerate}
  \item \textit{Khả năng áp dụng theo phiên bản kernel:} Dải phiên bản được trích xuất từ các mẫu như
        ``kernel versions prior to 5.12'' và so sánh với phiên bản kernel của máy chủ.
  \item \textit{Kiến trúc:} Các đề cập như ``x86\_64 only'' hoặc ``ARM-specific'' được khớp và
        các kiến trúc không phù hợp bị loại trừ.
  \item \textit{Yêu cầu quyền hạn:} Các CVE yêu cầu \texttt{CAP\_SYS\_ADMIN} hoặc các
        quyền hạn nâng cao khác bị loại trừ nếu container không có quyền hạn đó.
  \item \textit{Ràng buộc thiết bị:} Các khai thác yêu cầu phần cứng cụ thể (ví dụ:
        \texttt{/dev/kvm}, các thiết bị USB cụ thể) bị loại trừ nếu thiết bị không có sẵn.
\end{enumerate}

\subsection{Mã Giả Thuật Toán}
\label{subsec:appendix_pseudocode}

\begin{lstlisting}[float=htbp, language=Python, caption={Bộ phân tích CVE dựa trên quy tắc: thuật toán trích xuất syscall.}]
import re

def extract_high_risk_syscalls(cves, host_kernel, host_arch):
    """
    Extract the set of high-risk syscalls from CVE descriptions.

    Parameters:
        cves: List of CVE objects with a .desc field
        host_kernel: Host kernel version (tuple, e.g., (6, 8, 0))
        host_arch: Host architecture (str, e.g., "x86_64")

    Returns:
        S_blocked: Set of syscall names to block
    """
    S_blocked = set()

    for cve in cves:
        # Filter 1: Check kernel version applicability
        if not kernel_version_applicable(cve.desc, host_kernel):
            continue

        # Filter 2: Skip if architecture does not match
        if architecture_mismatch(cve.desc, host_arch):
            continue

        # Extract syscall mentions via regex
        found_syscalls = re.findall(syscall_regex, cve.desc)

        for syscall in found_syscalls:
            # Filter 3: Skip if special capability required
            #           but container does not have it
            if requires_special_cap(cve.desc) and \
               not container_has_cap():
                continue

            # Filter 4: Skip if specific device required
            #           but not available
            if needs_specific_device(cve.desc) and \
               not device_available():
                continue

            # Syscall passed all filters -> add to blocked set
            S_blocked.add(syscall)

    return S_blocked
\end{lstlisting}

\subsection{Ví Dụ Lọc Thực Tế}
\label{subsec:appendix_examples}

\textbf{Ví dụ~1: CVE-2022-0185~\cite{cve_2022_0185} (Linux Kernel --- Tràn bộ đệm fsconfig)}
\begin{itemize}
  \item \textit{Mô tả CVE:} ``A heap-based buffer overflow was found in the Linux kernel's
        legacy\_parse\_param function in the Filesystem Context functionality\ldots allows a
        local user with CAP\_SYS\_ADMIN to escalate privileges\ldots''
  \item \textit{Syscall Trích Xuất:} \texttt{ioctl}, \texttt{mount}
  \item \textit{Quyết Định Lọc:} Yêu cầu \texttt{CAP\_SYS\_ADMIN} $\Rightarrow$ giữ lại trong
        $S_{\text{blocked}}$ nếu container có quyền hạn đó; loại bỏ trong trường hợp ngược lại.
  \item \textit{Kết Quả:} Đối với các container không có đặc quyền, CVE này bị lọc ra. Đối với các container
        có \texttt{CAP\_SYS\_ADMIN}, \texttt{mount} được thêm vào $S_{\text{blocked}}$.
\end{itemize}

\textbf{Ví dụ~2: CVE-2021-22555~\cite{cve_2021_22555} (NetFilter --- tràn bộ đệm heap)}
\begin{itemize}
  \item \textit{Mô tả CVE:} ``A heap out-of-bounds write affecting Linux since
        v2.6.19-rc1 was discovered in net/netfilter/x\_tables.c\ldots can be exploited with
        setsockopt() to achieve local privilege escalation.''
  \item \textit{Syscall Trích Xuất:} \texttt{setsockopt}
  \item \textit{Điều Kiện Phiên Bản:} Ảnh hưởng đến kernel $< 5.12$; bị loại bỏ nếu kernel máy chủ là
        5.15+.
  \item \textit{Quyết Định Lọc:} Nếu kernel máy chủ là 5.15, CVE này bị bỏ qua và
        \texttt{setsockopt} \textit{không} được thêm vào $S_{\text{blocked}}$, tránh việc chặn không cần thiết
        các chức năng mạng hợp lệ.
\end{itemize}

Hai ví dụ này minh họa cách bộ phân tích CVE nhận thức ngữ cảnh tránh được cả kết quả dương tính giả
(chặn các syscall an toàn trong một số môi trường) và kết quả âm tính giả (bỏ sót các syscall nguy hiểm
khi áp dụng). Phương pháp tất định này đảm bảo rằng $S_{\text{blocked}}$ phản ánh chính xác
bề mặt tấn công CVE thực sự của mỗi môi trường container được triển khai.
